\section{Methodology}
\label{sec:methodology}

Our approach flows directly from the key insight: if Llama-2-7B layers vary in attention concentration, and high-entropy (diffuse) layers are functionally local, then a per-layer entropy score computed over a small calibration set can identify which layers are candidates for zero-shot SWA conversion. Every design choice below is motivated by this insight.

\subsection{Overview}

The pipeline comprises three stages:

\begin{enumerate}
\item \textbf{Entropy Scoring:} Compute per-layer attention entropy for all 32 layers using 100 calibration sequences.
\item \textbf{Layer Selection:} Rank layers by entropy; select the top-$k$ highest-entropy layers as SWA candidates.
\item \textbf{SWA Conversion:} Monkey-patch the selected layers with a sliding window attention mask; evaluate without fine-tuning.
\end{enumerate}

Stage~1 (entropy scoring) has been implemented and validated (h-e1). Stages~2 and~3 are fully designed and pending experimental execution (h-e2).

\subsection{Per-Layer Attention Entropy Scoring}

\subsubsection{Entropy Computation}

For each input sequence and each transformer layer, we extract the full attention weight tensor via \texttt{output\_attentions=True} in the HuggingFace forward pass (requires \texttt{attn\_implementation="eager"} for Llama-2-7B). Given the attention weight matrix $A \in \mathbb{R}^{h \times T \times T}$ for a layer with $h$ attention heads and sequence length $T$, we compute entropy as follows:

\begin{equation}
H_{\text{head}}(\text{head}, q) = -\sum_{k} A[\text{head}, q, k] \cdot \log\bigl(A[\text{head}, q, k] + \varepsilon\bigr)
\end{equation}

with $\varepsilon = 10^{-9}$ for numerical stability. Per-layer entropy is then:

\begin{equation}
H_{\text{layer}} = \underset{\text{heads}}{\text{mean}} \Bigl( \underset{q}{\text{mean}} \bigl( H_{\text{head}} \bigr) \Bigr)
\end{equation}

\textbf{Why head-mean pooling?} This is the critical design choice. Head-mean pooling captures the \textbf{layer-level average behavior} across all heads. Head-max pooling, by contrast, reports the single highest-entropy head, which is dominated by outlier heads that spread weight broadly regardless of the layer's typical behavior. We confirmed empirically that head-mean pooling yields Gini $= 0.681$ while head-max yields Gini $= 0.466$---a 32\% relative reduction in measured concentration when switching from head-mean to head-max---confirming that head-mean captures a more stable, layer-representative signal.

\subsubsection{Calibration Set Protocol}

We use the WikiText-103 validation split as the calibration source, following standard practice from quantization calibration (GPTQ, AWQ) which uses 100 sequences. We tokenize with the Llama-2-7B tokenizer (BPE, max\_length=2048) and chunk into contiguous non-overlapping sequences of 2048 tokens. Calibration subset A uses the first 100 sequences (indices 0--99). For stability analysis, subsets B and C use indices 100--199 and 200--299 respectively.

\textbf{Why 100 sequences?} This is sufficient for stable layer rankings (confirmed: Spearman $\rho \geq 0.8$ across all subset pairs) while requiring only $\approx 20$--30 GPU-minutes for 100 forward passes through Llama-2-7B on a single H100 (per calibration subset).

\textbf{Why WikiText-103?} It is the standard benchmark for our evaluation metric (perplexity) and represents the domain in which accuracy-preservation is measured. Using the same domain for calibration and evaluation minimizes distributional mismatch.

\subsubsection{Stability Validation}

We validate ranking stability using Spearman rank correlation $\rho$ between per-layer entropy vectors from different calibration subsets. The minimum $\rho$ across all subset pairs ($\rho_{\min} = \min(\rho_{AB}, \rho_{AC}, \rho_{BC})$) serves as the gate criterion: $\rho_{\min} \geq 0.8$ indicates that the entropy ranking is robust to calibration subset selection. This validation (Assumption A1) is a prerequisite for using entropy as a selection criterion.

\subsection{Layer Selection}

Given per-layer entropy scores $H \in \mathbb{R}^{32}$ from the calibration pass, we select the $k$ layers with the highest entropy as SWA candidates:

\begin{equation}
\text{selected\_layers} = \text{argsort}(H, \text{descending=True})[:k]
\end{equation}

\textbf{Why highest entropy?} High entropy indicates diffuse attention---the layer distributes weight broadly without sharp global token retrieval. These layers are not exploiting long-range global dependencies as strongly as low-entropy layers (which show sharp, concentrated attention on specific tokens). Replacing a high-entropy layer with SWA($w=512$) constrains it to attend within a local window---behavior closer to what it was already doing.

\textbf{Why $k = 4$?} Conservatively converting 12.5\% of layers (4 of 32) balances efficiency gain with risk of accuracy degradation. The remaining 28 full-attention layers continue to propagate global context through the residual stream, providing compensation for the 4 converted layers. We also investigate $k = 8$ (25\% conversion) as an upper bound in h-m2.

\subsection{Sliding Window Attention Conversion}

For each selected layer $l \in \text{selected\_layers}$, we monkey-patch the forward pass to apply an additive sliding window causal mask before the softmax:

\begin{lstlisting}[language=Python, basicstyle=\small\ttfamily, breaklines=true]
def swa_forward(self, hidden_states, attention_mask, ...):
    attn_weights = torch.matmul(query_states,
        key_states.transpose(-2,-1)) / math.sqrt(head_dim)
    seq_len = attn_weights.size(-1)
    swa_mask = make_sliding_window_causal_mask(
        seq_len, window=512, device=attn_weights.device)
    attn_weights = attn_weights + swa_mask  # 0.0 or -inf
    attn_weights = F.softmax(attn_weights, dim=-1)
    # ... continue standard attention
\end{lstlisting}

\textbf{Why additive float mask (0.0/$-\infty$)?} This follows HuggingFace's standard attention mask convention for Llama-2-7B (\texttt{attn\_implementation="eager"}), ensuring compatibility with the existing forward pass without requiring any structural modifications to the model. The mask is computed once and cached for efficiency.

\textbf{Why $w = 512$?} This matches the evaluation stride used for WikiText-103 perplexity computation (stride $= 512$), ensuring that the SWA window covers the full evaluation context for each token position. A window smaller than the evaluation stride would systematically exclude tokens that fall outside the window but within the stride context.

\subsection{Implementation Details}

\textbf{Model:} \texttt{meta-llama/Llama-2-7b-hf}, loaded with \texttt{attn\_implementation="eager"}, \texttt{torch\_dtype=torch.float16}, \texttt{device\_map="auto"}. Batch size $= 1$ (required for \texttt{output\_attentions=True} in float16 to avoid OOM). The attention tensor ($1 \times 32 \times 2048 \times 2048 \times$ float16) is deleted after each layer's entropy computation to avoid memory accumulation.

\textbf{Mask Validation:} Before evaluation, we execute a \texttt{verify\_swa\_mechanism()} check that confirms: (1) attended positions for each SWA layer match the expected sliding window pattern; (2) no off-by-one or causal mask interaction bugs are present.

\textbf{Evaluation Protocol for h-e2 (Pending):} WikiText-103 test split, stride $= 512$, max\_length $= 4096$. Perplexity computed as $\exp(\text{mean NLL})$ over the test set, consistent with standard HuggingFace perplexity benchmarks. Comparison conditions: (a) full-attention baseline ($k=0$), (b) entropy-guided $k=4$ SWA, (c) random-$k=4$ SWA (3 seeds, mean reported), (d) last-$k=4$ SWA (deepest 4 layers).

\subsection{Code Structure}

The implementation follows a clean module structure to enable reproducibility:

\begin{verbatim}
h-e1/code/
  data.py    -- load_wikitext103(), chunk_and_split()
  entropy.py -- compute_layer_entropy(), score_subset()
  report.py  -- compute_spearman(), save_results_json()
  run.py     -- end-to-end orchestration

h-e2/code/  [designed; execution pending]
  swa.py     -- make_sliding_window_causal_mask()
  eval.py    -- compute_perplexity(), compare_conditions()
  verify.py  -- verify_swa_mechanism()
  run.py     -- end-to-end orchestration
\end{verbatim}

All code is designed for single-GPU execution (1$\times$ H100), completing calibration for one subset (h-e1) in $\approx 20$--30 GPU-minutes; three subsets for full stability validation take approximately 1--1.5 GPU-hours. Evaluation (h-e2, pending) requires approximately 2--4 GPU-hours.
